\chapter{More Than Two Voices}

\section*{9.1}
A global description records interactions and causal ordering across roles, information no one local binary endpoint necessarily sees.

\section*{9.2}
Projection extracts from a global description the send/receive/branch obligations relevant to one role.

\section*{9.3}
It is a specification; projected participants can execute without a component that centrally schedules every event.

\section*{9.4}
Worker: receive `job:Nat` from broker, then send `result:Bool` to broker.

\section*{9.5}
Client: send `job:Nat` to broker, then receive `result:Bool` from broker.

\section*{9.6}
Let $A$ choose labels $l_1/l_2$ when talking to $B$, but require $C$ to send different next messages depending on that label without informing $C$. The local behavior for $C$ is underdetermined.

\section*{9.7}
Pairwise endpoint types can each look compatible while the combined causal ordering or branch knowledge cannot be realized globally.

\section*{9.8}
Exhaustion establishes equality of two finite trace sets for registered fixtures. A general theorem quantifies over a defined class of all well-formed/projectable global types.

\section*{9.9}
The unbranched choreography has exactly one global event sequence. Each event appears as complementary send/receive in the two projections, so synchronized composition reconstructs the same four-event sequence.

\section*{9.10}
For each global branch, any uninvolved role whose later behavior differs across branches must receive enough information to choose a unique local continuation. Violation makes projection ambiguous.

\section*{9.11}
Report the branch point, choosing role, dependent role, and the first pair of incompatible local continuations lacking a distinguishing message.

\section*{9.12}
Full credit requires an exact MPST source and theorem/hypothesis statement; citing the area generically is insufficient.
